% Copyright (C) 2014-2016 by Thomas Auzinger <thomas@auzinger.name>

\documentclass[draft,final]{vutinfth} % Remove option 'final' to obtain debug information.

% Load packages to allow in- and output of non-ASCII characters.
\usepackage{lmodern}        % Use an extension of the original Computer Modern font to minimize the use of bitmapped letters.
\usepackage[T1]{fontenc}    % Determines font encoding of the output. Font packages have to be included before this line.
\usepackage[utf8]{inputenc} % Determines encoding of the input. All input files have to use UTF8 encoding.

% Extended LaTeX functionality is enables by including packages with \usepackage{...}.
\usepackage{amsmath}    % Extended typesetting of mathematical expression.
\usepackage{amssymb}    % Provides a multitude of mathematical symbols.
\usepackage{mathtools}  % Further extensions of mathematical typesetting.
\usepackage{microtype}  % Small-scale typographic enhancements.
\usepackage[inline]{enumitem} % User control over the layout of lists (itemize, enumerate, description).
\usepackage{multirow}   % Allows table elements to span several rows.
\usepackage{booktabs}   % Improves the typesettings of tables.
\usepackage{subcaption} % Allows the use of subfigures and enables their referencing.
\usepackage[ruled,linesnumbered,algochapter]{algorithm2e} % Enables the writing of pseudo code.
\usepackage[usenames,dvipsnames,table]{xcolor} % Allows the definition and use of colors. This package has to be included before tikz.
\usepackage{nag}       % Issues warnings when best practices in writing LaTeX documents are violated.
\usepackage{todonotes} % Provides tooltip-like todo notes.
\usepackage{hyperref}  % Enables cross linking in the electronic document version. This package has to be included second to last.
\usepackage[acronym,toc]{glossaries} % Enables the generation of glossaries and lists fo acronyms. This package has to be included last.

% Define convenience functions to use the author name and the thesis title in the PDF document properties.
\newcommand{\authorname}{Rebeka Koszticsak} % The author name without titles.
\newcommand{\thesistitle}{Illustrative Thumbnails} % The title of the thesis. The English version should be used, if it exists.

% Set PDF document properties
\hypersetup{
    pdfpagelayout   = TwoPageRight,           % How the document is shown in PDF viewers (optional).
    linkbordercolor = {Melon},                % The color of the borders of boxes around crosslinks (optional).
    pdfauthor       = {\authorname},          % The author's name in the document properties (optional).
    pdftitle        = {\thesistitle},         % The document's title in the document properties (optional).
    pdfsubject      = {Subject},              % The document's subject in the document properties (optional).
    pdfkeywords     = {a, list, of, keywords} % The document's keywords in the document properties (optional).
}

\setpnumwidth{2.5em}        % Avoid overfull hboxes in the table of contents (see memoir manual).
\setsecnumdepth{subsection} % Enumerate subsections.

\nonzeroparskip             % Create space between paragraphs (optional).
\setlength{\parindent}{0pt} % Remove paragraph identation (optional).

\makeindex      % Use an optional index.
\makeglossaries % Use an optional glossary.
%\glstocfalse   % Remove the glossaries from the table of contents.

% Set persons with 4 arguments:
%  {title before name}{name}{title after name}{gender}
%  where both titles are optional (i.e. can be given as empty brackets {}).
\setauthor{}{\authorname}{}{female}
\setadvisor{Dr. techn.}{Manuela Waldner}{Msc.}{female}

% For bachelor and master theses:
%\setfirstassistant{Pretitle}{Forename Surname}{Posttitle}{male}
%\setsecondassistant{Pretitle}{Forename Surname}{Posttitle}{male}
%\setthirdassistant{Pretitle}{Forename Surname}{Posttitle}{male}

% For dissertations:
\setfirstreviewer{Pretitle}{Forename Surname}{Posttitle}{male}
\setsecondreviewer{Pretitle}{Forename Surname}{Posttitle}{male}

% For dissertations at the PhD School and optionally for dissertations:
\setsecondadvisor{Pretitle}{Forename Surname}{Posttitle}{male} % Comment to remove.

% Required data.
\setaddress{Address}
\setregnumber{1325492}
\setdate{01}{01}{2001} % Set date with 3 arguments: {day}{month}{year}.
\settitle{\thesistitle}{Illustrative Thumbnails} % Sets English and German version of the title (both can be English or German).
%\setsubtitle{Optional Subtitle of the Thesis}{Optionaler Untertitel der Arbeit} % Sets English and German version of the subtitle (both can be English or German).

% Select the thesis type: bachelor / master / doctor / phd-school.
% Bachelor:
\setthesis{bachelor}
%
% Master:
%\setthesis{master}
%\setmasterdegree{dipl.} % dipl. / rer.nat. / rer.soc.oec. / master
%
% Doctor:
%\setthesis{doctor}
%\setdoctordegree{rer.soc.oec.}% rer.nat. / techn. / rer.soc.oec.
%
% Doctor at the PhD School
%\setthesis{phd-school} % Deactivate non-English title pages (see below)

% For bachelor and master:
\setcurriculum{Media Informatics and Visual Computing}{Medieninformatik und Visual Computing} % Sets the English and German name of the curriculum.

% For dissertations at the PhD School:
\setfirstreviewerdata{Affiliation, Country}
\setsecondreviewerdata{Affiliation, Country}


\begin{document}

\frontmatter % Switches to roman numbering.
% The structure of the thesis has to conform to
%  http://www.informatik.tuwien.ac.at/dekanat

\addtitlepage{naustrian} % German title page (not for dissertations at the PhD School).
\addtitlepage{english} % English title page.
\addstatementpage

\begin{danksagung*}
\todo{Ihr Text hier.}
\end{danksagung*}

\begin{acknowledgements*}
\todo{Enter your text here.}
\end{acknowledgements*}

\begin{kurzfassung}
\todo{Ihr Text hier.}
\end{kurzfassung}

\begin{abstract}
Thumbnails are used to display the screens when switching between them on the computer
and on mobile devices. These images make it easier to recognize the opened applications,
and help to find the needed window quicker. Thumbnails display however only a screenshot
of the windows, so they get potentially confusing if there are more opened windows or if the
same application is opened multiple times. Depending on the resolution of the display, the
screenshot size decreases as the number of opened windows increases. Furthermore,
within the same application (like MS Office World) the screenshots are similar in appearance
(eg.: white paper and tool bar), but the important text is not readable.
There are several approaches that filter the important areas of the images to make editting
less obvious or enhance the main region. In this bachelor thesis an application
is implemented that uses these methods on the screencaptured
images. The less important
areas of the screenshots are cut off, and the thumbnails show only important information,
which makes them more illustrative and easier to fulfill their purpose.
\end{abstract}

% Select the language of the thesis, e.g., english or naustrian.
\selectlanguage{english}

% Add a table of contents (toc).
\tableofcontents % Starred version, i.e., \tableofcontents*, removes the self-entry.

% Switch to arabic numbering and start the enumeration of chapters in the table of content.
\mainmatter

\chapter{Introduction}
\todo{Enter your text here.}

\chapter{Related Work}
There are several image processing algorithms, which can be helpful by creating illustrative thumbnails.
The main difference between them is, if they consider the fact that the input images always are screenshots.
Therefore, this section is divided into two parts.
The first one discusses algorithms with UI processing segments.
Algorithms, invented for retrieve visual data from common images, are examinated in the second section. 
According to the information visualization method, like combining the most important parts in form of collages or simple resizing, two classes are divided.
In the following some of these methods are compared.

\section{Processing as UI}
Considering that the input is a screenshot, there is a high chance that common UI elements are shown on the screen.
Exceptions are only the cases where graphics, images, videos or application containing these, e.g. video games, gallery program, video player are captured.
The applications like that tend to hide all UI elements, "fullscreen" mode, or redefine them, e.g. game menu.\par
Labeling the image parts as content and non-content, the metadata about the UI elements can be helpful.
Forras uses already existing accessibility APIs to segment UI and non-UI data. 
Matching the metadata with the screen content provides a fast and robust result about the location of any kind of UI content.
There are howerer several disadvantages, that need to be take into account when using such APIs.
The range and the granularity of the support is often not wide enough.
The use of an accessibilty API does not make sure, that every UI element will be recognized, because some metadata is not reachable or it will be ignored by the application.\par
Because of that Forras Prefab works with its own prototypes. 
In the database models and  prototypes of common UI elements are saved.
The (components of) UI elements has to be matched with the prototypes in the data base, and after that the predefined metadata can be accessed.
Since the Prefab system is able to split complex widgets into their base elements, the database does not need to be unnecessary big, but it is still able to cover the most common UI elements.
In the case of special or rare UI widgets, like in a video game application or elements of a not not widely used software, the system fails.
If a widget is not saved in the database, there is no way, that it will be recognized.\par
Forras Sikuli offers a solution not for the incompleteness of the database, but for the granularity issues too.
It uses its own templates just like the Prefab system, but only in case of small icons and widgets.
Since in case of larger objects template matching would be too expensive, after the processing of a training pattern, the Sikuli system is able to create new object models too.
Although this feature is used in the application for other purpose, i.e. reduce the matching cost, it can solve the problem of database and of the granularity.
Sikuli makes it possible to expand the database and to make the database entry afterwards more detailed.\par
But in case of accessibility APIs not only the availability of metadata causes possible problems.
It has also no information about the actual visibilty of the UI elements.
One window or rather one widget can hide an other one, some content can be out of the range of the borders of the screen etc.    
Forras is based on the Prefab system, but it builds a hierarchical tree from the widgets. 
The content is found at the leafs, and the parents are the widget, where the child is built in.
With the help of this tree the order of the UI elements gets clear, and so the misleading information can be eliminated.\par
After the labeling of the UI elements correctly, they can be processed as needed.
They can cut off as whole or processed according to the information content.
Forras invented an algorithm to eliminate objects of a picture and fill their place with the texture of the object behind them using the z-buffer information.
The tree mentioned above works as a z-buffer in this case.
With this cut off algorithm unnecessary widgets can be easily eliminated and more important elements of the UI can so get better visible.\par
According to the definition of "illustrative" of this thesis the UI elements of a screen has to get cut so the segmentation of the UI elements is not needed.
Although the Prefab and Sikuli systems can be helpful for labeling the images as UI and content.
But theses approaches has the drawback of the usage of the database and the slow template matching. 
Since there is no need to know, which widgets are on the screen, and processing the actual content can already cause performance issues, the use these methods would overcomplicate the application without providing noteworthy advantages.

\section{Processing as common picture}
There are several information retrieving methods for processing images with any content.
Furthermore in case of making any kind of image more illustrative, there are important tasks like interesting point recognition, Region of Interest (ROI) selection, image or feature composition etc. which they can handle easily.
According to the manner of the input data there are two different groups of these algorithms, which need to be discussed in the following sections.
The first category works with more than one picture in the same time.
Its strength is to choose singe features which represents the most the whole input data. 
In exchange it is likely that no input image will be recognizable on the result.
On the other hand there are the methods in the second group, which take only one picture for input and process it as one unit.
So the result is similar to the input data, but therefore the chance is high that not only the unimportant but the areas with high information ratio are damaged
At the very and of this section some approaches are presented, which can help the work of both of the above mentioned groups.

\subsection{Collage creating methods} 
A collage is an assembled image, containing parts of a bunch of input images and being representative for the whole input data.
There are two cases by creating thumbnails more illustrative where such methods can be helpful.
On the one hand the actual information of a screenshot image concentrates only in few regions of the picture.
Many parts, for example UI elements, space between the content etc, can be ignored. 
The other way around the content can be retrieved in form of ROIs, and afterwards they can be combined arbitrary.
On the other hand using collage creator algorithms thumbnails for desktop switching can be easily generated, where the input are screenshots of the open application of the desktop instead of some ROIs of one screen.\par
For a representative collage the most important task is to choose the best images, which information content covers the whole input data.
Forras takes the parameters representativeness, importance costs, transition cost and object sensitivity into account.
Representativeness means being interesting in this case.
A picture tends to have high representativeness value, if there are many special textures on it, and if it is not similar to the rest of the data (so that no image is chosen twice).
Importance cost evaluates and collect the ROIs of the input.
While transition cost stands for the smooth transition between every two images.
At last, the parameter object sensitivity holds the results of object recognition, and it helps in the arrangement, that the object has a reasonable placement.\par
Forras concentrates however only the first two parameters above.
It clusters the images, according their source and the time, and measures the quality.
Thank for the clustering by the choice of the final images it is clear,  which images are the same or have similar content.
This feature, accordingly modified, can be useful by sorting ROIs of a screenshot, i.e.: text content, image content etc. or of the running applications of the desktop, i.e.: textprocessing, gallery application etc.
The parameter quality summarizes the value of the results of the following calculations: blurriness, compression, contrast and color balance.
Since in this case only screenshots, thus computer generated pictures, can be the input, these measurements invented for camera data would provide less meaningful results than the algorithm above.\par 
Having the best ROIs for the collage the last task is to merge them into one output picture.
For this purpose Forras uses a method called ROI packing.
First the central point of every ROI needs to get chosen and every pixel on the canvas needs to get assigned to one of the using the k-Means algorithm.
After that the ROIs can be placed on the area calculated for them.
To fill the place between the ROIs they have to grow, keeping the aspect ratio, until they overlap. 
Then every ROI is shifted to the middle of its area.
This method can be repeated until no ROI grows anymore.
To fill the white areas, neighboring ROIs are allowed to cover them, and so fill the whole image.\par
Collage methods are very useful to represent a large image dataset only in a small place.
They work with numerous input data, and taking the most important parts of them create a new image, that is not similar to any before.
That is why they are more useful for making a thumbnail for a desktop but not for one application.
Even though taking only the most important ROIs of one screen it would be possible to create an image more illustrative that any other, since the important content is the largest and the best readable.
Cutting one screen apart and arranging some parts of them willingly would confuse the user, and even more time would be needed to recognize the screen.

\subsection{Resampling methods}
Treating the input image as one unit, has a big advantage against the approaches mentioned before.
Even after modification the result will be highly similar to the input.
Resampling means, that some parts of the image will be eliminated, but all remaining areas will have the same proportion to each other, not like by the collage algorithms.
The image remains recognizable because of the same look, even though the most important areas are less readable.\par 
To select areas, which needs to kept same, Forras uses several attention models, which defines Attention Objects (AO).
AOs usually are objects from the real world, which catch the human eye, because of their familiarity, shape, color etc.
With three values AOs can be easily parametrized, these are ROI, Attention Value (AV) and Minimal Perceptible Size (MPS).
The attention models fit the AOs into the context.
Forras works with three different attention model at the same time, with saliency, face and text attention.
Having the importance value of every pixel the most important area can be detected.\par 
The algorithm above after carefully calculation chooses only one area, which contains as many AOs as possible.
So some possibly important AOs have to be ignored and cut off.
With Feature awared Texturing described in Forras this do not have to be the case.
The algorithm expects an input image and a feature mask.
A grid is generated, which lies on the input image.
This grid can be modified in an optional shape, but the gridpoints on the feature mask are not allowed to change their proportion to each other.
In that way the picture elements between the AOs fill the new shape, but the AOs get barely distorted\par 
A detailed importance is essential by creating thumbnails, since a screen usually contain more information and ROIs than a common picture. 
But the algorithms described above has aspects, like face recognition and grid calculation, which overcomplicate the calculations, causing performance issues without reaching better quality.
A face on a computer screen is not so common as on usual photos, and it is not so important too than text for example.
With a grid the input image can get reshaped in any other form, and in addition no important part needs to get damaged.
But it is meant to form the input in completely other shape, not for resize it according to aspect ratio.
Some aspects however, like definition of AOs, can be useful polishing the actual algorithm, this will be discussed in the future work section.

\subsection{Feature combining methods}

\chapter{Methodology}



% Remove following line for the final thesis.
%%% intro.tex
%% Copyright (C) 2014-2016 by Thomas Auzinger <thomas@auzinger.name>
%
% This work may be distributed and/or modified under the
% conditions of the LaTeX Project Public License, either version 1.3
% of this license or (at your option) any later version.
% The latest version of this license is in
%   http://www.latex-project.org/lppl.txt
% and version 1.3 or later is part of all distributions of LaTeX
% version 2005/12/01 or later.
%
% This work has the LPPL maintenance status `maintained'.
%
% The Current Maintainer of this work is Thomas Auzinger.
%
% This work consists of the files vutinfth.dtx and vutinfth.ins
% and the derived file vutinfth.cls.
% This work also consists of the file intro.tex.


\newacronym{ctan}{CTAN}{Comprehensive TeX Archive Network}
\newacronym{faq}{FAQ}{Frequently Asked Questions}
\newacronym{pdf}{PDF}{Portable Document Format}
\newacronym{svn}{SVN}{Subversion}
\newacronym{wysiwyg}{WYSIWYG}{What You See Is What You Get}

\newglossaryentry{texteditor}
{
  name={editor},
  description={A text editor is a type of program used for editing plain text files.}
}

\chapter{Introduction to \LaTeX}

Since \LaTeX\ is widely used in academia and industry, there exists a plethora of freely accessible introductions to the language.
Reading through the guide at \url{https://en.wikibooks.org/wiki/LaTeX} serves as a comprehensive overview for most of the functionality and is highly recommended before starting with a thesis in \LaTeX.

\section{Installation}

A full \LaTeX\ distribution\index{distribution} consists of not only of the binaries that convert the source files to the typeset documents, but also of a wide range of packages and their documentation.
Depending on the operating system, different implementations are available as shown in Table~\ref{tab:distrib}.
\textbf{Due to the large amount of packages that are in everyday use and due to their high interdependence, it is paramount to keep the installed distribution\index{distribution} up to date.}
Otherwise, obscure errors and tedious debugging ensue.

\begin{table}
  \centering
  \begin{tabular}{cccc}
    \toprule
    Distribution & Unix         & Windows      & MacOS        \\
    \midrule
    TeX Live     & \textbf{yes} & yes          & (yes)        \\
    MacTeX       & no           & no           & \textbf{yes} \\
    MikTeX       & no           & \textbf{yes} & no           \\
    \bottomrule
  \end{tabular}
  \caption{\TeX/\LaTeX\ distributions for different operating systems. Recomended choice in \textbf{bold}.}
  \label{tab:distrib} % \label has to be placed AFTER \caption to produce correct cross-references.
\end{table}

\section{Editors}

A multitude of \TeX\ \glspl{texteditor} are available differing in their editing models, their supported operating systems and their feature sets.
A comprehensive overview of \glspl{texteditor} can be found at the Wikipedia page  \url{https://en.wikipedia.org/wiki/Comparison_of_TeX_editors}.
TeXstudio (\url{http://texstudio.sourceforge.net/}) is recommended.
Most editors support the scrolling the typeset preview document to a location in the source document by \verb|Ctrl| clicking the location in the source document.

\section{Compilation}

Modern editors usually provide the compilation programs to generate \gls{pdf} documents and for most \LaTeX\ source files, this is sufficient.
More advanced \LaTeX\ functionality, such as glossaries and bibliographies, needs additional compilation steps, however.
It is also possible that errors in the compilation process invalidate intermediate files and force subsequent compilation runs to fail.
It is advisable to delete intermediate files (\verb|.aux|, \verb|.bbl|, etc.), if errors occur and persist.
All files that are not generated by the user are automatically regenerated.
To compile the current document, the steps as shown in Table~\ref{tab:compile} have to be taken.


\begin{table}
  \centering
  \begin{tabular}{rl}
    \toprule
    & Description \\
    \midrule
    1 & Scan for refs, toc/lof/lot/loa items and cites \\
    2 & Build the bibliography     \\
    3 & Link refs and build the toc/lof/lot/loa \\
    4 & Link the bibliography \\
    5 & Build the glossary \\
    6 & Build the acronyms \\
    7 & Build the index \\
    8 & Link the glossary, acronyms, and the index \\
    9 & Link the bookmarks \\
    \midrule
    & Command \\
    \midrule
    1 & \verb|pdflatex.exe  example| \\
    2 & \verb|bibtex.exe    example| \\
    3 & \verb|pdflatex.exe  example| \\
    4 & \verb|pdflatex.exe  example| \\
    5 & \verb|makeindex.exe -t example.glg -s example.ist| \\
      & \verb|              -o example.gls example.glo| \\
    6 & \verb|makeindex.exe -t example.alg -s example.ist| \\
      & \verb|              -o example.acr example.acn| \\
    7 & \verb|makeindex.exe -t example.ilg -o example.ind example.idx| \\
    8 & \verb|pdflatex.exe  example| \\
    9 & \verb|pdflatex.exe  example| \\
    \bottomrule
  \end{tabular}
  \caption{Compilation steps for this document. The following abbreviations were used: table of contents (toc), list of figures (lof), list of tables (lot), list of algorithms (loa).}
  \label{tab:compile} % \label has to be placed AFTER \caption to produce correct cross-references.
\end{table}


\section{Basic Functionality}

In this section, various examples are given of the fundamental building blocks used in a thesis.
Many \LaTeX\ commands have a rich set of options that can be supplied as optional arguments.
The documentation of each command should be consulted to get an impression of the full spectrum of its functionality.

\subsection{Floats}

'Two main categories of page elements can be differentiated in the usual \LaTeX\ workflow: \textit{(i)} the main stream of text and \textit{(ii)} floating containers that are positioned at convenient positions throughout the document.
In most cases, tables, plots, and images are put into such containers since they are usually positioned at the top or bottom of pages.
These are realized jtk dgm,fsds by the two environments \verb|figure| and \verb|table|, which also provide functionality for cross-referencing (see Table~\ref{tab:intro} and Figure~\ref{fig:intro}) and the generation of corresponding entries in the list of figures and the list of tables.
Note that these environments solely act as containers and can be assigned arbitrary content.

\subsection{Tables}

A table in \LaTeX\ is created by using a \verb|tabular| environment or any of its extensions, e.g., \verb|tabularx|.
The commands \verb|\multirow| and \verb|\multicolumn| allow table elements to span multiple rows and columns.

\begin{table}[h] % placement specifier
  \centering
  \begin{tabular}{lll}
    \toprule
    \multicolumn{2}{c}{Position} \\
    \cmidrule{1-2} % partial horizontal rule
    Group & Abbrev & Name \\
    \midrule
    Goalkeeper & GK & Paul Robinson \\
    \midrule
    \multirow{4}{*}{Defenders} & LB & Lucus Radebe \\
                               & DC & Michael Duburry \\
                               & DC & Dominic Matteo \\
                               & RB & Didier Domi \\
    \midrule
    \multirow{3}{*}{Midfielders} & MC & David Batty \\
                                 & MC & Eirik Bakke \\
                                 & MC & Jody Morris \\
    \midrule
    Forward & FW & Jamie McMaster \\
    \midrule
    \multirow{2}{*}{Strikers} & ST & Alan Smith \\
                              & ST & Mark Viduka \\
    \bottomrule
  \end{tabular}
  \caption{Adapted example from the \LaTeX guide at \url{https://en.wikibooks.org/wiki/LaTeX/Tables}. This example uses rules specific to the \texttt{booktabs} package and employs the multi-row functionality of the \texttt{multirow} package.}
  \label{tab:intro} % \label has to be placed AFTER \caption to produce correct cross-references.
\end{table}

\subsection{Images}

An image is added to a document via the \verb|\includegraphics| command as shown in Figure~\ref{fig:intro}.
The \verb|\subcaption| command can be used to reference subfigures, such as Figure~\ref{fig:intro:full width} and~\ref{fig:intro:half width}.

\begin{figure}[h]
  \centering
  \begin{subfigure}[b]{0.45\columnwidth}
    \centering
    \includegraphics[width=\textwidth]{TU_INF_Logo_gray}
    \subcaption{The header logo at text width.}
    \label{fig:intro:full width}
  \end{subfigure}
  \begin{subfigure}[b]{0.45\columnwidth}
    \centering
    \includegraphics[width=0.5\textwidth]{TU_INF_Logo_gray}
    \subcaption{The header logo at half the text width.}
    \label{fig:intro:half width}
  \end{subfigure}
  \caption{The header logo at different sizes.}
  \label{fig:intro} % \label has to be placed AFTER \caption (or \subcaption) to produce correct cross-references.
\end{figure}

\subsection{Mathematical Expressions}

One of the original motivation to create the \TeX\ system was the need for mathematical typesetting.
To this day, \LaTeX\ is the preferred system to write math-heavy documents and a wide variety of functions aids the author in this task.
A mathematical expression can be inserted inline as $\sum_{n=1}^{\infty} \frac{1}{n^2} = \frac{\pi^2}{6}$ outside of the text stream as \[ \sum_{n=1}^{\infty} \frac{1}{n^2} = \frac{\pi^2}{6} \] or as numbered equation with
\begin{equation}
\sum_{n=1}^{\infty} \frac{1}{n^2} = \frac{\pi^2}{6}.
\end{equation}

\subsection{Pseudo Code}

The presentation of algorithms can be achieved with various packages; the most popular are \verb|algorithmic|, \verb|algorithm2e|, \verb|algorithmicx|, or \verb|algpseudocode|.
An overview is given at \url{https://tex.stackexchange.com/questions/229355}.
An example of the use of the \verb|alogrithm2e| package is given with Algorithm~\ref{alg:gauss-seidel}.

\begin{algorithm}
  \SetKw{BreakFor}{break for}
  \KwIn{A scalar~$\epsilon$, a matrix $\mathbf{A} = (a_{ij})$, a vector $\vec{b}$, and an initial vector $\vec{x}^{(0)}$}
  \KwOut{$\vec{x}^{(n)}$ with $\mathbf{A} \vec{x}^{(n)} \approx \vec{b}$}
  \For{$k\leftarrow 1$ \KwTo maximum iterations}
  {
     \For{$i\leftarrow 1$ \KwTo $n$}
     {
        $x_i^{(k)} = \frac{1}{a_{ii}} \left(b_i-\sum_{j<i} a_{ij} x_j^{(k)} - \sum_{j>i} a_{ij} x_j^{(k-1)} \right)$\;
     }
     \If{$\lvert\vec{x}^{(k)}-\vec{x}^{(k-1)}\rvert < \epsilon$}
     {\BreakFor\;}
  }
  \Return{$\vec{x}^{(k)}$\;}
  \caption{Gauss-Seidel}
  \label{alg:gauss-seidel} % \label has to be placed AFTER \caption to produce correct cross-references.
\end{algorithm}

\section{Bibliography}

The referencing of prior work is a fundamental requirement of academic writing and well supported by \LaTeX.
The \textsc{Bib}\TeX\ reference management software is the most commonly used system for this purpose.
Using the \verb|\cite| command, it is possible to reference entries in a \verb|.bib| file out of the text stream, e.g., as~\cite{Turing1936}.
The generation of the formatted bibliography needs a separate execution of \verb|bibtex.exe| (see Table~\ref{tab:compile}).

\section{Table of Contents}

The table of contents is automatically built by successive runs of the compilation, e.g., of \verb|pdflatex.exe|.
The command \verb|\setsecnumdepth| allows the specification of the depth of the table of contents and additional entries can be added to the table of contents using \verb|\addcontentsline|.
The starred versions of the sectioning commands, i.e., \verb|\chapter*|, \verb|\section*|, etc., remove the corresponding entry from the table of contents.

\section{Acronyms / Glossary / Index}

The list of acronyms, the glossary, and the index need to be built with a separate execution of \verb|makeindex| (see Table~\ref{tab:compile}).
Acronyms have to be specified with \verb|\newacronym| while glossary entries use \verb|\newglossaryentry|.
Both are then used in the document content with one of the variants of \verb|\gls|, such as \verb|\Gls|, \verb|\glspl|, or \verb|\Glspl|.
Index items are simply generated by placing \verb|\index|\marg{entry} next to all the words that correspond to the index entry \meta{entry}.
Note that many enhancements exist for these functionalities and the documentation of the \verb|makeindex| and the \verb|glossaries| packages should be consulted.

\section{Tips}

Since \TeX\ and its successors do not employ a \gls{wysiwyg} editing scheme, several guidelines improve the readability of the source content:
\begin{itemize}
\item Each sentence in the source text should start with a new line.
      This helps not only the user navigation through the text, but also enables revision control systems (e.g. \gls{svn}, Git) to show the exact changes authored by different users.
      Paragraphs are separated by one (or more) empty lines.
\item Environments, which are defined by a matching pair of \verb|\begin{name}| and \verb|\end{name}|, can be indented by whitespace to show their hierarchical structure.
\item In most cases, the explicit use of whitespace (e.g. by adding \verb|\hspace{4em}| or \verb|\vspace{1.5cm}|) violates typographic guidelines and rules.
      Explicit formatting should only be employed as a last resort and, most likely, better ways to achieve the desired layout can be found by a quick web search.
\item The use of bold or italic text is generally not supported by typographic considerations and the semantically meaningful \verb|\emph{|\texttt{$\dots$}\verb|}| should be used.
\end{itemize}

The predominant application of the \LaTeX\ system is the generation of \gls{pdf} files via the \textsc{Pdf}\LaTeX\ binaries.
In the current version of \textsc{Pdf}\LaTeX, it is possible that absolute file paths and user account names are embedded in the final \gls{pdf} document.
While this poses only a minor security issue for all documents, it is highly problematic for double blind reviews.
The process shown in Table~\ref{tab:ps2pdf} can be employed to strip all private information from the final \gls{pdf} document.

\begin{table}[h]
  \centering
  \begin{tabular}{rl}
  \toprule
  & Command \\
  \midrule
  1 & Rename the \gls{pdf} document \verb|final.pdf| to \verb|final.ps|. \\
  2 & Execute the following command: \\
    & \verb|ps2pdf -dPDFSETTINGS#/prepress ^| \\
    & \verb| -dCompatibilityLevel#1.4 ^| \\
    & \verb| -dAutoFilterColorImages#false ^| \\
    & \verb| -dAutoFilterGrayImages#false ^| \\
    & \verb| -dColorImageFilter#/FlateEncode ^| \\
    & \verb| -dGrayImageFilter#/FlateEncode ^| \\
    & \verb| -dMonoImageFilter#/FlateEncode ^| \\
    & \verb| -dDownsampleColorImages#false ^| \\
    & \verb| -dDownsampleGrayImages#false ^| \\
    & \verb| final.ps final.pdf| \\
  \bottomrule
  \end{tabular}

  On Unix-based systems, replace \verb|#| with \verb|=| and \verb|^| with \verb|\|.
  \caption{Anonymization of \gls{pdf} documents.}
  \label{tab:ps2pdf}
\end{table}

\section{Resources}

\subsection{Useful Links}

In the following, a listing of useful web resources is given.
\begin{description}
\item[\url{https://en.wikibooks.org/wiki/LaTeX}] An extensive wiki-based guide to \LaTeX.
\item[\url{http://www.tex.ac.uk/faq}] A (huge) set of \gls{faq} about \TeX\ and \LaTeX.
\item[\url{https://tex.stackexchange.com/}] The definitive user forum for non-trivial \LaTeX-related questions and answers.
\end{description}

\subsection[Comprehensive TeX Archive Network]{\gls{ctan}}

The \gls{ctan} is the official repository for all \TeX\ related material.
It can be accessed via \url{https://www.ctan.org/} and hosts (among other things) a huge variety of packages that provide extended functionality for \TeX\ and its successors.
Note that most packages contain \gls{pdf} documentation that can be directly accessed via \gls{ctan}.

In the following, a short, non-exhaustive list of relevant \gls{ctan}-hosted packages is given together with their relative path.
\begin{description}[itemsep=0ex]
\item[\href{https://www.ctan.org/pkg/algorithm2e}{algorithm2e}] Functionality for writing pseudo code.
\item[\href{https://www.ctan.org/pkg/amsmath}{amsmath}] Enhanced functionality for typesetting mathematical expressions.
\item[\href{https://www.ctan.org/pkg/amsfonts}{amssymb}] Provides a multitude of mathematical symbols.
\item[\href{https://www.ctan.org/pkg/booktabs}{booktabs}] Improved typesetting of tables.
\item[\href{https://www.ctan.org/pkg/enumitem}{enumitem}] Control over the layout of lists (\verb|itemize|, \verb|enumerate|, \verb|description|).
\item[\href{https://www.ctan.org/pkg/fontenc}{fontenc}] Determines font encoding of the output.
\item[\href{https://www.ctan.org/pkg/glossaries}{glossaries}] Create glossaries and list of acronyms.
\item[\href{https://www.ctan.org/pkg/graphicx}{graphicx}] Insert images into the document.
\item[\href{https://www.ctan.org/pkg/inputenc}{inputenc}] Determines encoding of the input.
\item[\href{https://www.ctan.org/pkg/l2tabu}{l2tabu}] A description of bad practices when using \LaTeX.
\item[\href{https://www.ctan.org/pkg/mathtools}{mathtools}] Further extension of mathematical typesetting.
\item[\href{https://www.ctan.org/pkg/memoir}{memoir}] The document class on upon which the \verb|vutinfth| document class is based.
\item[\href{https://www.ctan.org/pkg/multirow}{multirow}] Allows table elements to span several rows.
\item[\href{https://www.ctan.org/pkg/pgfplots}{pgfplots}] Function plot drawings.
\item[\href{https://www.ctan.org/pkg/pgf}{pgf/TikZ}] Creating graphics inside \LaTeX\ documents.
\item[\href{https://www.ctan.org/pkg/subcaption}{subcaption}] Allows the use of subfigures and enables their referencing.
\item[\href{https://www.ctan.org/tex-archive/info/symbols/comprehensive/}{symbols/comprehensive}] A listing of around 5000 symbols that can be used with \LaTeX.
\item[\href{https://www.ctan.org/pkg/voss-mathmode}{voss-mathmode}] A comprehensive overview of typesetting mathematics in \LaTeX.
\item[\href{https://www.ctan.org/pkg/xcolor}{xcolor}] Allows the definition and use of colors.
\end{description} % A short introduction to LaTeX.

\backmatter

% Use an optional list of figures.
\listoffigures % Starred version, i.e., \listoffigures*, removes the toc entry.

% Use an optional list of tables.
\cleardoublepage % Start list of tables on the next empty right hand page.
\listoftables % Starred version, i.e., \listoftables*, removes the toc entry.

% Use an optional list of alogrithms.
\listofalgorithms
\addcontentsline{toc}{chapter}{List of Algorithms}

% Add an index.
\printindex

% Add a glossary.
\printglossaries

% Add a bibliography.
\bibliographystyle{alpha}
\bibliography{intro}

\end{document}